	
\begin{chineseabstract}
本文档为应用于 14-bit 高精度逐次逼近型模数转换器（SAR ADC）的模拟前端（AFE）
及全差分电容数模转换器（CDAC）的仿真验证报告。针对工业现场微弱信号的采集需求，
本报告基于 TSMC $180\,\mathrm{nm}$ CMOS 工艺，在典型 $3.3\,\mathrm{V}$ 单电源供电及多种 
PVT（工艺、电压、温度）角下，对跨阻放大器（TIA）、单端转差分电路（S2D）、
全差分放大器（FDA）以及可编程增益放大器（PGA）进行了较为系统的交直流、瞬态及噪声协同仿真。

仿真结果表明，AFE 整链在不同增益配置下的有效带宽基本达到 $700\,\mathrm{kHz}$，
全链路折算至输入端的总积分噪声约为 $45\,\mathrm{\mu V_{rms}}$，满量程大信号输出总谐波失真
（THD）优于 $-88\,\mathrm{dB}$。针对后级等效约 $40\,\mathrm{pF}$ 的容性负载，PGA 
在 $800\,\mathrm{ns}$ 内较好地实现了目标精度的建立。此外，蒙特卡洛（Monte Carlo）
分析初步验证了版图匹配策略对失调电压的抑制作用，以及连续时间共模反馈（CMFB）环路的有效性。
综合仿真数据表明，本电路在动态范围、信噪比（SNR）等指标上基本满足 14-bit 系统的设计预期
；但需指出，目前的仿真数据多基于原理图级及部分寄生估计，实际流片后，封装寄生与全局布线等
非理想因素可能对系统最终的噪声与建立性能产生一定影响。



\chinesekeyword{时域电磁散射，时域积分方程，时间步进算法，后时不稳定性，时域平面波算法}
\end{chineseabstract}

